\begin{helper}
For every natural number \(N\) and every vector
\(c\in\mathbb R^{\operatorname{Fin}(N)}\), there is a vector
\(x\in\mathbb R^{\operatorname{Fin}(N)}\) and a permutation
\(\sigma\) of \(\operatorname{Fin}(N)\) such that
\[
x_i=c_{\sigma(i)}\quad\text{for all }i
\]
and
\[
i\le j\implies x_j\le x_i .
\]
In other words, every finite real vector admits a nonincreasing
permuted ordering, with ties broken arbitrarily.
\end{helper}
\begin{proof}
We argue by an explicit finite selection procedure.  Let
\(U=\operatorname{Fin}(N)\).  If \(N=0\), then \(U\) is empty; take the unique
empty vector and the unique empty permutation.  The asserted equations and
order condition are vacuous.

Assume now that \(N>0\).  Since \(U\) is finite and the set of already chosen
coordinates will always be finite, the following recursive construction is
well defined.  At step \(0\), choose an index \(\sigma(0)\in U\) on which
\(c\) attains its maximum over all of \(U\).  At step \(r+1<N\), after
\(\sigma(0),\ldots,\sigma(r)\) have been chosen distinctly, choose
\(\sigma(r+1)\) on which \(c\) attains its maximum over the nonempty finite
set
\[
U\setminus\{\sigma(0),\ldots,\sigma(r)\}.
\]
If several indices have the same value, choose any one of them.  This tie
choice is harmless because equal values may be ordered in either way.

By construction the selected indices are all distinct.  There are exactly
\(N\) steps and \(U\) has exactly \(N\) elements, so every element of \(U\) is
selected once.  Thus the assignment \(r\mapsto \sigma(r)\), after identifying
the ordinal positions \(0,\ldots,N-1\) with \(\operatorname{Fin}(N)\), is a
bijection of \(U\), hence a permutation.

Define \(x_r=c_{\sigma(r)}\).  If \(r\le s\), then when \(\sigma(r)\) was
chosen, the later element \(\sigma(s)\) had not yet been chosen and therefore
belonged to the remaining candidate set.  Since \(\sigma(r)\) was chosen to
maximize \(c\) on that candidate set,
\[
c_{\sigma(s)}\le c_{\sigma(r)}.
\]
Equivalently, \(x_s\le x_r\).  This proves the required nonincreasing order,
and the defining equation \(x_i=c_{\sigma(i)}\) proves that \(x\) is a
permutation of \(c\).
\end{proof}
